\section{Objetivos}

\subsection{Objetivo General}
Observar y evaluar experimentalmente el comportamiento del rendimiento de un transformador monofásico de laboratorio, separando el efecto de las pérdidas en el hierro y las pérdidas en los arrollamientos, y verificar las condiciones operativas y de montaje conforme a las recomendaciones de las normativas vigentes.

\subsection{Objetivos Específicos}
\begin{enumerate}
    \item Determinar el rendimiento experimental del transformador monofásico a partir de la relación entre la potencia activa de salida y de entrada, validando la descomposición de las pérdidas totales en pérdidas en el núcleo y pérdidas en los devanados.
    \item Evaluar el impacto de las pérdidas sin carga ($P_{\text{Fe}}$) obtenidas del ensayo de circuito abierto y de las pérdidas bajo carga ($I_2^2 R_{\text{eq2}}$) calculadas a partir del ensayo de cortocircuito sobre la curva de eficiencia.
    \item Obtener experimentalmente la curva de rendimiento $\eta$ en función de la corriente de carga secundaria $I_2$ para un factor de potencia unitario (carga resistiva pura).
    \item Analizar el comportamiento del rendimiento bajo régimen de carga inductiva ($\cos\varphi < 1$), determinando la influencia del factor de potencia sobre la eficiencia global del transformador.
    \item Localizar y comprobar el punto de rendimiento máximo ($I_2^* = \sqrt{P_{\text{Fe}}/R_{\text{eq2}}}$), verificando la condición teórica de igualdad entre las pérdidas fijas en el hierro y las pérdidas variables en el cobre.
    \item Estimar la potencia nominal de la unidad ensayada a partir de las coordenadas del punto de máxima eficiencia ($S^* = V_2^* I_2^*$) y compararla con los datos de placa de la máquina.
    \item Aplicar el método de la resistencia para medir la elevación térmica del devanado secundario al término del régimen de plena carga, corrigiendo la resistencia óhmica a la temperatura de operación real.
    \item Emplear los lineamientos de la norma COVENIN 3172-95 (punto 5.7) y el estándar ANSI/IEEE C57.12.91-2001 para la determinación normalizada del aumento de temperatura en los devanados, y contrastar los límites térmicos admisibles con COVENIN 536-94.
\end{enumerate}
